  \item \field{Feasibility for Real-World Use}

    \textbf{Prototype feasibility.} The Round-1 goal is a working prototype, not a
    finished medical product. The hardware is realistic because it uses commodity parts,
    but the hard work is integration, testing Thai-language prompts, and learning how
    real users respond in home routines. Comparable smart mirrors have already been
    built~[6, 7].

    \textbf{Privacy and trust by design.} A camera and microphone in the home demand
    trust. The prototype therefore starts with on-device processing where possible,
    minimal data retention, no continuous cloud streaming of faces, and clear
    family/caregiver controls~[14].

    \textbf{Adoption strategy.} To avoid the complexity and poor support that drive
    assistive-technology abandonment~[12, 13], the pilot uses co-design, shallow
    menus, caregiver setup, and a light support pathway.

    \textbf{Cost and who pays.} The indicative prototype cost is about 8{,}500\,THB
    before assembly, casing refinement, support, and maintenance. That is affordable
    for a prototype but still not cheap for many users, so likely early funders are
    families with support, disability foundations, vocational-rehabilitation programmes,
    employers under the disability quota, and public agencies.

    \textbf{Path to real use.} The practical next step is a small home pilot: verify
    that prompts are understood, caregivers can configure routines, privacy settings are
    trusted, and the device survives ordinary use. Partnerships with DEP, disability
    foundations, community health services, vocational-rehabilitation centres, employers,
    and caregiver networks would be needed before any wider rollout.
